\documentclass[10pt,a4paper]{amsart}
\usepackage[margin=19mm]{geometry}
\usepackage{amsmath,amssymb,mathtools}
\usepackage[scr=rsfs,cal=euler]{mathalfa}
\usepackage{xfrac}
\usepackage[unicode,hidelinks,bookmarks=false]{hyperref}
\newcommand{\ie}{i.e.\ }
\newcommand{\cf}{cf.\ }
\newcommand{\resp}{respectively\ }
\newcommand{\Subsection}{\subsection}
\providecommand{\nobreakdash}{\nobreak\hbox{-}}
\setlength{\emergencystretch}{2em}
\multlinegap0pt
\usepackage{bm}
\usepackage{bbm}
\usepackage{dsfont}
\usepackage{xspace}
\usepackage{cancel}
\usepackage{mathtools}
\usepackage{amsthm}
\makeatletter
\@ifundefined{Thm}{\newtheorem{Thm}{Thm}}{}
\@ifundefined{Claim}{\newtheorem{Claim}[Thm]{Claim}}{}
\makeatother
\def\R{\mathbb{R}}
\title[Canonical torus action on symplectic singularities]{Canonical torus action on symplectic singularities}
\author{Yoshinori Namikawa, Yuji Odaka}
\date{Source: arXiv:2503.15791v4 (CC BY 4.0)}
\begin{document}
\maketitle

\noindent\emph{Source and licence.} Canonical torus action on symplectic singularities. Version arXiv:2503.15791v4 (CC BY 4.0), distributed under the Creative Commons Attribution 4.0 International licence, as recorded in the arXiv licence metadata for this exact version and shown on the abstract page https://arxiv.org/abs/2503.15791v4. Full rights notice: licenses/arxiv-2503.15791-CC-BY-4.0.txt.\par\smallskip

\noindent\textit{Editorial scope.} Counted unit: the Claim whose source label is \texttt{claim:completevf3} in the article (source0.tex lines 5389--5392, source label \texttt{claim:completevf3}), delivered together with the article's own complete original proof of it (source0.tex lines 5393--5456, one \texttt{proof} environment of 64 lines). No mathematics is written, repaired or re-derived: statement and proof are the article's own text line for line. Required context retained uncounted: none. Every definition, display and label of those lines is retained unchanged. Omitted from this package and not redistributed: 1--5388, 5457--6190. Nothing inside a retained excerpt is omitted or altered: every retained body is delivered exactly as the article's lines are written, so no omission receipt of any kind accompanies this package. Citations: the retained excerpts carry no citation command at all, so no bibliography entry of the article is transcribed and no reference is redistributed. Reference adaptation: none was needed and none was made; every reference inside the delivered unit resolves inside it, so no unresolved reference or citation is produced. Printed slips kept literal: no printed word, symbol, hypothesis or number of the counted statement or of its proof was altered. Layout and class: the article's own class is \texttt{amsart} with packages including amsmath, amssymb, amsthm, amscd, epsfig, enumerate, enumitem, color, bbm, todonotes, version, geometry; the wrapper is the lane's pinned \texttt{amsart} editorial wrapper at 10pt on A4, which restates the theorem-like environments the retained text needs and adds the article's own macro definitions that the retained text needs (\texttt{\textbackslash R}), with extra wrapper packages (bm, bbm, dsfont, xspace, cancel, mathtools). The delivered numbering is the unit's own. Assets: no retained span contains a figure environment, a graphics-inclusion command, a PDF-page inclusion, a table or a TikZ picture; no image file is copied into this package, the package declares no asset dependency and no image file is redistributed with it. Licence: version arXiv:2503.15791v4 of this article is distributed under the Creative Commons Attribution 4.0 International licence, as recorded in the arXiv licence metadata for this exact version and shown on the abstract page https://arxiv.org/abs/2503.15791v4. A full rights notice is delivered at licenses/arxiv-2503.15791-CC-BY-4.0.txt. Characters: every retained excerpt is pure ASCII, so the delivered engine, XeLaTeX with its default Latin Modern text font, reports no missing character.
\begin{Claim}\label{claim:completevf3}
For any $d,e,f\in \R$, 
$d\xi_I+e\xi_J+f\xi_K$ on $C_x(X)^{sm}$ is complete. 
\end{Claim}

\begin{proof}
For the general case of Claim \ref{claim:completevf3} 
(possibly with $f \neq 0$), we employ the well-known isomorphism 
$\mathfrak{sp}(1) \simeq \mathfrak{o}(3)$ as follows
to make the arguments more intuitively 
understandable. 
We recall its basic as follows. 
Write the quaternion as $\mathbb{H} = \R \oplus \R I \oplus \R J \oplus \R K$, 
define $Sp(1) = \{q \in \mathbb{H}\: \vert \: q\bar{q} = 1\}$ and consider its 
the adjoint action of $Sp(1)$ on $\mathfrak{sp}(1) :=\R I \oplus \R J \oplus \R K$ as $q\colon x \mapsto qxq^{-1}$ for $x \in \mathfrak{sp}(1)$. 
Through the natural identification $\R^3 = \R I \oplus \R J \oplus \R K$, 
the adjoint action induces a $2:1$ 
covering $Sp(1) \twoheadrightarrow SO(3)$ of Lie groups and an isomorphism 
$\mathfrak{sp}(1) \cong \mathfrak{o}(3)$ of Lie algebras.

%identifying 
%$I, J, K$ with the generators of rotations in 
%$\mathbb{R}^3$ to make the arguments more intuitively 
%understandable. 
The proof then reduces to the previous case by the 
following trick. 

Below, we show 
that for any $d,e,f\in \R$, 
$dI+eJ+fK \in \R^3$ can be written as 
$({\rm exp}(aI+bJ))\cdot (cI)$ with the exponential map 
${\rm exp}\colon \mathfrak{sp}(1)\simeq \mathfrak{o}(3)
\to SO(3)$. Here, this $\cdot$ means the standard 
action. 
Indeed, $a,b,c$ can be taken using $d,e,f$ explicitly as follows. 
First, $c$ must represent the norm of the vector, i.e., 
$c = \sqrt{d^2 + e^2 + f^2}$. Set $d':=\frac{d}{c}, e':=\frac{e}{c}, 
f':=\frac{f}{c}$. 
We also want to also introduce $\theta$ as angle required to rotate the $x$-axis 
to the direction of $\R(d, e, f)$ with 
the rotation axis as $\R(a,b,0)$. 
We can define them explicitly: 
let $(a',b')\in \R^2$ as 
$a'^2+b'^2=1$ and $(a',b',0)\perp (1-
d',e',f')$. Explicitly, 
$a':=\frac{e'}{\sqrt{(1-d')^2+e'^2}}, 
b':=\frac{d'-1}{\sqrt{(1-d')^2+e'^2}}$. 
Then, since the distance between 
$(a',b')$ and $(1,0,0)$ is the same as 
that between $(d',e',f')$ and $(a',b')$, 
there is $\theta\in [0,2)$ 
so that angle $\pi\theta (\in [0, 2\pi\theta))$ rotation 
along the axis $\R_{\ge 0}(a',b')$ brings 
$(1,0,0)$ to $(d',e',f')$. Accordingly, 
we set 
$a:= a' \theta, \quad b:=b' \theta$. 
Then, the element $dI+eJ+fK$ is realized 
as the image of $cI$ 
under the adjoint action of $\exp(aI+bJ)$. 
Thus, 
\begin{align*}
\phi(dI+eJ+fK)&=\phi(({\rm exp}(aI+bJ))(cI))\\
&=({\rm exp}(aI+bJ))_*\phi(cI)
\end{align*}
is still complete on $C_x(X)^{sm}$ because 
it is simply the pushforward of a complete vector field by a (global) 
diffeomorphism. We finish the proof of 
Claim \ref{claim:completevf3}. 
\end{proof}

\end{document}
